\section{Continuous products, envelopes and domination}
\label{sec:hessian-continuous-products}\index{Hessian measure!continuous potentials}

Uniform admissible approximation allows us to define the mixed Hessian
products. Their local bounds retain the derivatives of the Hermitian
background; closedness of the background is not required.

\Needspace{13\baselineskip}
\begin{proposition}\label{hessian:continuous-products}
For continuous admissible $u_1,\ldots,u_k$, $1\leq k\leq m$, the current
\[
 (\omega+\ddc u_1)\wedge\cdots\wedge(\omega+\ddc u_k)
                                          \wedge\omega^{n-m}
\]
is well defined as the limit of uniform smooth admissible approximations.
It is positive, local, and continuous when its potentials converge
uniformly. Its local masses are bounded in terms of uniform bounds on
the potentials. In particular $H_m(u)$ is a positive Radon measure.
\end{proposition}
\begin{proof}
\Cref{hessian:smooth-proper-existence,hessian:continuous-smooth-approximation} provide the smooth approximations without
using any continuous Hessian product or domination input.
For a smooth admissible $u$ put $\sigma=\omega+\ddc u$.
Since $\mathrm d\sigma=\mathrm d\omega$, the coefficients of
$\ddc(\chi\sigma^{r-1}\wedge\omega^{n-r})$ are sums of smooth test
factors paired with powers $\sigma^q$, $q\leq r-1$, where $\chi$ is a
fixed cutoff. The absolute-form lemma and integration by parts give
the Chern--Levine--Nirenberg bound by induction on $r\leq m$:
\[
 \int_K\sigma^r\wedge\omega^{n-r}
                     \leq C_{K,U,\omega}(1+\|u\|_{L^\infty(U)})^r.
\]
At the induction step write $\sigma^r=\omega\wedge\sigma^{r-1}
+\ddc u\wedge\sigma^{r-1}$, integrate the second term against $\chi$,
and move $\ddc$ to $\chi\sigma^{r-1}\omega^{n-r}$.
The resulting powers have degrees $r-1,r-2,r-3$, with negative powers absent, and the induction
applies on nested slightly larger compact sets. Mixed masses reduce to
the single potential average, by positivity of the expansion of its
power. Positive mixed currents have total variation controlled by these
trace masses.

To prove independence, telescope the difference between two smooth
mixed products whose potentials have a common uniform bound. Each term
paired with a fixed smooth complementary test is
\[
 \int\ddc(u_a-v_a)\wedge T\wedge\eta
                         =\int(u_a-v_a)\ddc(T\wedge\eta),
\]
where $T$ is a product of the remaining admissible factors and
$\omega^{n-m}$. The identity $\mathrm d(\omega+\ddc u)=\mathrm d\omega$
and its $\dc$ analogue show that derivatives of $T$ introduce only
fixed metric derivatives; every term has at most $k-1<m$ admissible
factors. The absolute mixed-form bound and the mass estimate just
proved bound the variation of $\ddc(T\wedge\eta)$ uniformly.
Consequently the pairing difference is at most
$C_\eta\max_a\|u_a-v_a\|_\infty$.
The same construction with an arbitrary complementary smooth test
also defines the signed products of fewer than $m$ factors without the
factor $\omega^{n-m}$: their variation is bounded by the absolute-form
lemma. These products justify the metric-derivative terms above.
This proves that all approximations have the same limit and that the
products are continuous in the uniform topology. Positivity passes
from the smooth mixed products. Using a test supported in an open set
on which two limiting potentials coincide proves locality by the same
estimate on a smaller neighborhood of its support.
\end{proof}

Write $H_m(u)$ for the measure with every potential equal to $u$.
Locality and uniform continuity of these products give comparison on a
contact set.

\begin{lemma}\label{hessian:contact-comparison}
If $u\leq v$ are continuous admissible functions and the remaining
$m-1$ factors are fixed continuous admissible functions, then on
$\{u=v\}$ the mixed measure with first factor $\omega+\ddc u$ is at
most the one with first factor $\omega+\ddc v$.
In particular $\mathbf1_{\{u=v\}}H_m(u)\leq
\mathbf1_{\{u=v\}}H_m(v)$.
\end{lemma}
\begin{proof}
Write $M(w)$ for the mixed measure with first potential $w$.
Positivity and locality give, for $t>0$,
\[
 M(\max(v,u+t))\geq
 \mathbf1_{\{v>u+t\}}M(v)+\mathbf1_{\{u\geq v\}}M(u).
\]
The first indicators increase to $\mathbf1_{\{v>u\}}$ as $t\downarrow0$,
and the potentials on the left converge uniformly to $v$.
Integrate against nonnegative continuous tests and pass to the limit.
Subtract the first restricted measure to obtain the claimed contact
inequality. Replacing one factor at a time proves the pure-power claim
and the analogous contact inequalities for all mixed powers.
\end{proof}

For a continuous real obstacle $h$, define $P(h)$ to be the supremum of
continuous admissible functions bounded above by $h$. Constants make the
family nonempty. The penalty construction also proves that this
supremum is continuous and its Hessian measure lies on the contact set.
\index{Envelope!continuous Hessian obstacle}

\begin{lemma}\label{hessian:continuous-obstacle-envelope}
The function $P(h)$ is continuous and admissible, and $H_m(P(h))$ is
supported on $\{P(h)=h\}$. Moreover
$\|P(h_1)-P(h_2)\|_\infty\leq\|h_1-h_2\|_\infty$.
\end{lemma}
\begin{proof}
For smooth $h$, use the penalty solutions $w_N$ in \cref{hessian:continuous-smooth-approximation}. Its maximum argument applies to every continuous admissible
competitor $v\leq h$, with an error depending on $h$ and $N$, and not
on $v$. Taking their supremum gives
\[
 P(h)-E_N-2M_N/N\leq w_N\leq P(h),
\]
where the error tends to zero. Thus $P(h)$ is their uniform limit and
is continuous and admissible. On every compact subset of
$\{P(h)<h\}$ the penalty densities
$b e^{N(w_N-h)}$ tend uniformly to zero. Product continuity implies
the contact support assertion.
Adding constants to competitors gives the stated Lipschitz bound.
Approximate a continuous $h$ uniformly by smooth obstacles. Their
envelopes converge uniformly to $P(h)$, proving continuity and
admissibility. Any compact subset of $\{P(h)<h\}$ also lies outside the
contact sets of all sufficiently close smooth obstacles; product
continuity gives the support assertion for $h$.
\end{proof}

The obstacle also bounds discontinuous competitors. Their upper-test
admissibility suffices for the penalty estimate, so allowing them does
not change the continuous envelope of a smooth obstacle.

\begin{proposition}\label{hessian:bounded-envelope}
For every smooth real obstacle $h$, the envelope $P(h)$ equals the
pointwise supremum of all bounded admissible functions below $h$.
\end{proposition}
\begin{proof}
The solutions $w_N$ of \cref{hessian:bounded-penalty-squeeze} are
continuous admissible competitors, so $w_N\leq P(h)$. For each
bounded admissible $u\leq h$, the same lemma gives
\[
 u\leq w_N+E_N+2M_N/N\leq P(h)+E_N+2M_N/N.
\]
Let $N\to\infty$ for this fixed obstacle and fixed $b$.
Thus $P(h)$ dominates every bounded competitor. Since $P(h)$ itself
is a bounded competitor, it equals their supremum.
\end{proof}

The trace test in \eqref{eq:hessian-continuous-admissibility} gives
$L_\omega u\geq-n$, where
$L_\omega u=n\ddc u\wedge\omega^{n-1}/\omega^n$.
We use the classical normalized Green estimate
\begin{equation}\label{eq:hessian-normalized-l1}
 \int_X(\sup_Xu-u)\omega^n\leq C_\omega
\end{equation}
for continuous admissible $u$. The stationary probability of $L_\omega$
has a smooth positive density equivalent to $\omega^n$. The positive
measure $L_\omega u+n$ has stationary mass $n$, so the Green-kernel
estimate gives this bound. It is a linear elliptic input, and does not
assert compactness for a nonlinear Hessian measure.

To use envelopes in domination, we must first exclude a potential with
zero total Hessian mass.

\begin{lemma}\label{hessian:hessian-nonvanishing}
Every continuous admissible $u$ satisfies $\int_X H_m(u)>0$.
\end{lemma}
\begin{proof}
Suppose this integral is zero. For $s>0$ put
$w=\max(u,\min_Xu+s)$ and $U=\{u<\min_Xu+s\}$.
The open set $U$ is nonempty, $0\leq w-u\leq s$, and $w=u$ outside it.
Set $\sigma_u=\omega+\ddc u$, $\sigma_w=\omega+\ddc w$, and
$\mu_q=\sigma_u^q\wedge\sigma_w^{m-q}\wedge\omega^{n-m}$.
The contact lemma gives $\mathbf1_{X\setminus U}\mu_q\geq
\mathbf1_{X\setminus U}\mu_m$, for $0\leq q\leq m$.
For $q<m$ the telescoping identity is
\[
 \mu_q-\mu_m=\ddc(w-u)\wedge T_q,\qquad
 T_q=\sigma_u^q\wedge
       \sum_{r=0}^{m-q-1}\sigma_w^r\wedge\sigma_u^{m-q-1-r}
                                      \wedge\omega^{n-m}.
\]
Integration by parts on $X$ therefore gives
\[
 \int_U(\mu_q-\mu_m)\leq\int_X(\mu_q-\mu_m)
                                =\int_X(w-u)\ddc T_q.
\]
This identity follows first for smooth approximants and then by uniform
product continuity. The derivatives of $T_q$ contain only fixed metric
derivatives and products of at most $m-1$ admissible factors, so the
absolute-form lemma makes $\ddc T_q$ a measure of uniformly controlled
variation. Uniform convergence also justifies its continuous weight.
On $U$, locality gives $\sigma_w=\omega$ and thus
\[
 |\ddc T_q|\leq C\sum_{r=0}^{m-1}\mu_r.
\]
The constant depends only on the metric and $m,n$, and not on $u$ or $s$.
Consequently
\[
 \sum_{q=0}^m\int_U\mu_q
 \leq(m+1)\int_U H_m(u)+Cms\sum_{r=0}^{m-1}\int_U\mu_r.
\]
For sufficiently small $s$, absorption and $H_m(u)=0$ force all these
nonnegative integrals to vanish. But $\mu_0=\omega^n$ on the nonempty
open set $U$, a contradiction.
\end{proof}

The envelope of a stretched difference now converts vanishing of a
restricted measure into an order relation. The linear bound
\eqref{eq:hessian-normalized-l1} prevents these envelopes from tending to
$-\infty$ on an open set.

\begin{proposition}\label{hessian:zero-hessian-domination}
If continuous admissible $u,v$ satisfy
$\mathbf1_{\{u<v\}}H_m(u)=0$, then $u\geq v$.
\end{proposition}
\begin{proof}
Replace $v$ by $\max(u,v)$, so $u\leq v$.
For $b>1$ put $p_b=P(bu-(b-1)v)$ and
$z_b=b^{-1}p_b+(1-b^{-1})v\leq u$.
On the contact set $C_b=\{p_b=bu-(b-1)v\}$ one has $z_b=u$;
the contact inequality and polarized positivity give
\[
 \mathbf1_{C_b}b^{-m}H_m(p_b)
                 \leq\mathbf1_{C_b}H_m(z_b)
                 \leq\mathbf1_{C_b}H_m(u).
\]
The envelope measure is supported on $C_b$ and has positive mass.
It vanishes on $C_b\cap\{u<v\}$ by the hypothesis, so some point of
$C_b\cap\{u=v\}$ exists. Hence
$\sup_Xp_b\geq\inf_Xv$, while $p_b\leq u$.
The estimate \eqref{eq:hessian-normalized-l1} is consequently
uniform for $p_b$. On the open set $E_a=\{u<v-a\}$, for $a>0$, one has
$p_b\leq v-ba$, and therefore
\[
 C_\omega\geq\int_X(\sup_Xp_b-p_b)\omega^n
                     \geq(ba-\operatorname{osc}_Xv)\int_{E_a}\omega^n.
\]
Letting $b\to\infty$ gives zero volume for $E_a$. Since it is open it
is empty, and $a\downarrow0$ gives $u\geq v$.
\end{proof}

\begin{proposition}\label{hessian:strict-hessian-domination}
If continuous admissible $u,v$ satisfy
\[
 \mathbf1_{\{u<v\}}H_m(u)\leq
                    c\mathbf1_{\{u<v\}}H_m(v),\qquad 0\leq c<1,
\]
then $u\geq v$.
\end{proposition}
\begin{proof}
Choose $b>1$ so that $(1-b^{-1})^m>c$, and put
$p=P(bu-(b-1)v)$ and $z=b^{-1}p+(1-b^{-1})v\leq u$.
On $C=\{p=bu-(b-1)v\}$ the contact comparison gives
\[
 \mathbf1_C\bigl(b^{-m}H_m(p)+(1-b^{-1})^mH_m(v)\bigr)
                                  \leq\mathbf1_C H_m(u).
\]
Restricting to $\{u<v\}$ and using the hypothesis shows that
$H_m(p)$ vanishes on $C\cap\{u<v\}$. Its support lies in $C$, where
$p<v$ is equivalent to $u<v$. Thus $H_m(p)$ vanishes on $\{p<v\}$.
The preceding domination proposition gives $p\geq v$.
Since $p\leq bu-(b-1)v$, it follows that $u\geq v$.
\end{proof}

Strict domination yields the comparison needed for the proper equations.
\index{Domination principle!continuous Hessian measures}

\begin{lemma}\label{hessian:hessian-exponential-order}
For continuous admissible $u,v$ and $a>0$, if
\[
 e^{-av}H_m(v)\geq e^{-au}H_m(u),
\]
then $u\geq v$.
\end{lemma}
\begin{proof}
Fix $\delta>0$. On $\{u<v-\delta\}$, the assumed inequality gives
\[
 H_m(u)\leq e^{a(u-v)}H_m(v)
       \leq e^{-a\delta}H_m(v-\delta).
\]
The constant $e^{-a\delta}$ is strictly smaller than one, so \cref{hessian:strict-hessian-domination} applied to $u$ and $v-\delta$ gives $u\geq v-\delta$. Letting $\delta\downarrow0$ proves the assertion.
\end{proof}
